\section{The lcm square map and cutoff return}
\label{sec:lcm}

Define the lcm coefficient
\begin{equation}
 \Gamma_I(q)=
 \sum_{\substack{d_1,d_2\in I\\{}[d_1,d_2]=q}}
 \mu(d_1)\mu(d_2).
 \label{eq:Gamma-definition}
\end{equation}

\begin{theorem}[Exact lcm transform]
\label{thm:lcm-transform}
For every $k\ge1$,
\begin{equation}
 \boxed{\ \beta_I(k)^2=\sum_{q\mid k}\Gamma_I(q).\ }
 \label{eq:lcm-transform}
\end{equation}
Consequently,
\begin{equation}
 \Gamma_I(q)=\sum_{a\mid q}\mu(q/a)\beta_I(a)^2.
 \label{eq:Gamma-inversion}
\end{equation}
Moreover,
\begin{equation}
 \supp\Gamma_I\subset\{q\le V^2:q\ \text{squarefree}\},
 \qquad
 |\Gamma_I(q)|\le3^{\omega(q)},
 \label{eq:Gamma-support-bound}
\end{equation}
and
\begin{equation}
 \sum_q|\Gamma_I(q)|\le V^2.
 \label{eq:Gamma-l1}
\end{equation}
\end{theorem}

\begin{proof}
Expand the square in \eqref{eq:beta-I}.  The conditions
$d_1\mid k$ and $d_2\mid k$ are equivalent to $[d_1,d_2]\mid k$;
grouping by that least common multiple proves \eqref{eq:lcm-transform}.
Equation \eqref{eq:Gamma-inversion} is ordinary M\"obius inversion.

Only squarefree $d_i$ contribute.  Hence their lcm is squarefree and at
most $V^2$.  For a squarefree $q$, each prime divisor can lie in
$d_1$ only, $d_2$ only, or both, giving at most $3^{\omega(q)}$
ordered pairs.  Finally, summing absolute values before grouping by the
lcm gives at most the number of ordered pairs $d_1,d_2\in I$, which is
at most $V^2$.
\end{proof}

The sign structure becomes more transparent after separating the gcd.

\begin{proposition}[Pairwise-coprime lcm coordinates]
\label{prop:gab}
One has the exact formula
\begin{equation}
 \Gamma_I(q)=
 \sum_{\substack{gab=q\\
 g,a,b\ \mathrm{squarefree}\\
 (g,a)=(g,b)=(a,b)=1\\
 D_0<ga\le V,\ D_0<gb\le V}}
 \mu(a)\mu(b),
 \label{eq:Gamma-gab}
\end{equation}
The squarefree condition is part of the summation range.
\end{proposition}

\begin{proof}
For a nonzero pair in \eqref{eq:Gamma-definition}, put
$g=(d_1,d_2)$, $a=d_1/g$, and $b=d_2/g$.  Squarefreeness makes
$g,a,b$ pairwise coprime, $[d_1,d_2]=gab$, and
\[
 \mu(d_1)\mu(d_2)
 =\mu(ga)\mu(gb)=\mu(g)^2\mu(a)\mu(b)=\mu(a)\mu(b).
\]
The construction is reversible.
\end{proof}

Since $V^2\le R$ by \eqref{eq:RV}, the lcm label created by squaring
never leaves the finite Ramanujan cutoff:
\begin{equation}
 \boxed{\quad q=[d_1,d_2]\le V^2\le R.\quad}
 \label{eq:cutoff-return}
\end{equation}
This is an algebraic support statement, not a cancellation estimate.
In particular, the diagonal pair $d_1=d_2=q$ contributes
\begin{equation}
 \Gamma_I(q)=\mu(q)^2\one_{q\in I}+\Gamma_I^{\ne}(q),
 \label{eq:Gamma-divisor-diagonal}
\end{equation}
and neither term on the right is automatically small.

For later use, summing \eqref{eq:lcm-transform} against a test function
$F$ gives the exact row assembly
\begin{equation}
 \sum_k\beta_I(k)^2F(k)
 =\sum_{q\le R}\Gamma_I(q)\sum_jF(qj).
 \label{eq:lcm-row-assembly}
\end{equation}
Thus the same-$k$ quadraticization has a small row label $q\le R$.
The $d$-first quadraticization of \cref{sec:opened-d} instead has the
large row label $\ell d>R$.  These are complementary normal forms, not
two notations for the same modulus.
